\section{Physics Geometry: Device Region Labeling}

\subsection{Files}

\begin{itemize}
  \item \texttt{src/physics/geometry/regions.py}
  \item \texttt{scripts/geometry/build\_device\_regions.py}
  \item \texttt{tests/test\_device\_regions.py}
\end{itemize}

\subsection{Physical Significance}

Device region labeling converts the calibrated geometry into a discrete physical
partition of the channel. The label map is device-level: it is defined once for
the chip geometry and then reused for every droplet frame. This prevents the
occupancy module from re-inventing spatial regions on the fly.

The labels are:

\begin{center}
\begin{tabular}{@{}cl@{}}
\toprule
ID & Physical region \\
\midrule
0 & background/unassigned \\
1 & inlet \\
2 & outlet \\
3 & left branch \\
4 & right branch \\
5 & upper junction \\
6 & lower junction \\
\bottomrule
\end{tabular}
\end{center}

\subsection{Governing Geometry}

The physical width of the channel is:

\[
W = 100~\mu\mathrm{m} = 25~\mathrm{px}.
\]

The region label map is constructed by assigning pixels to the closest
centerline branch or configured junction region while respecting the fixed
channel geometry. The region labels serve as the spatial basis functions for
occupancy:

\[
\Omega =
\Omega_{\mathrm{inlet}}
\cup
\Omega_{\mathrm{outlet}}
\cup
\Omega_{\mathrm{left}}
\cup
\Omega_{\mathrm{right}}
\cup
\Omega_{\mathrm{upper~junction}}
\cup
\Omega_{\mathrm{lower~junction}}.
\]

\subsection{Implementation Details}

The label builder:

\begin{itemize}
  \item reads the device geometry and calibration;
  \item uses preconfigured junction centers and physical junction sizes;
  \item builds integer labels with no nonzero labels outside the channel;
  \item validates that all physical regions are nonempty;
  \item validates connectivity between branches and junctions;
  \item creates a region overlay for manual inspection.
\end{itemize}

\subsection{Outputs}

\begin{verbatim}
data/geometry/asymmetric_loop_h100/region_labels.npy
data/geometry/asymmetric_loop_h100/region_metadata.json
data/geometry/asymmetric_loop_h100/region_overlay.png
\end{verbatim}

\begin{figure}[H]
  \centering
  \includegraphics[width=0.95\linewidth,height=0.75\textheight,keepaspectratio]{../data/geometry/asymmetric_loop_h100/region_overlay.png}
  \caption{Device-level physical region overlay. This image shows the precomputed region labels used by droplet occupancy. Regions are not recomputed by the occupancy or validation visualizers.}
  \label{fig:region-overlay}
\end{figure}

\begin{figure}[H]
  \centering
  \includegraphics[width=0.95\linewidth,height=0.75\textheight,keepaspectratio]{../outputs/geometry/asymmetric_loop_h100/centerline_regions_frame_0.png}
  \caption{Historical centerline-region diagnostic for frame 0. It is retained as a development figure; the production occupancy path uses the device-level labels in Figure~\ref{fig:region-overlay}.}
  \label{fig:centerline-regions-frame0}
\end{figure}

\subsection{Execution}

\begin{verbatim}
.\.venv\Scripts\python.exe -m scripts.geometry.build_device_regions \
  --experiment configs\experiments\video_2.yml \
  --frame-index 0 \
  --overwrite
\end{verbatim}
